% Originally: content/week-09.tex, \subsection{Length of a curve} --- promoted to a section

\section{Length of a curve}
\label{sec:length_of_curve}

% Source: Corsin Nick/Class Notes/Week 9.pdf, p. 8
\begin{definition}[Length of a $C^1$-curve]
\label{def:length_c1_curve}
The \newterm{length} (\germanterm{Länge}) of a curve or path $\gamma \in
C^1([a,b],\mathbb{R}^n)$ is given by
\[L(\gamma) := \int_a^b \lvert\gamma'(t)\rvert\,dt.\]
\end{definition}

\begin{definition}[Length of a $C^0$-curve]
\label{def:length_c0_curve}
For a $C^0$-curve $c \in C^0([a,b],\mathbb{R}^n)$, the length is defined as
\[L(c) := \sup\left\{\sum_{i=1}^{n-1} \lvert c(t_{i+1})-c(t_i)\rvert
  \ :\ a = t_1 < \cdots < t_n = b\right\},\]
the supremum over all lengths of inscribed polygons.
\end{definition}

% Sonnet 5 (Medium) --- FIG-W09-04
\begin{center}
\begin{tikzpicture}[>=Latex, scale=1.0]
  \draw[thick] plot[smooth] coordinates
    {(0,0) (0.6,0.9) (1.3,0.5) (1.9,0.15) (2.5,0.75) (3.1,0.5) (3.7,1.3)};
  \node[font=\scriptsize] at (1.9,1.15) {$\operatorname{Im}(c)$};
  \draw[blue, thick]
    (0,0) -- (0.55,0.95) -- (1.35,0.55) -- (2.0,0.1) -- (2.55,0.7) -- (3.15,0.45) -- (3.7,1.25);
  \foreach \x/\y/\lbl/\pos in {0/0/{c(t_1)}/below,0.55/0.95/{c(t_2)}/above,1.35/0.55/{c(t_3)}/above,2.0/0.1/{c(t_4)}/below,2.55/0.7/{c(t_5)}/{above left},3.15/0.45/{c(t_6)}/below,3.7/1.25/{c(t_7)}/right} {
    \fill (\x,\y) circle (1.3pt);
    \node[font=\scriptsize, \pos] at (\x,\y) {$\lbl$};
  }
\end{tikzpicture}
\end{center}

\begin{remark}
For a $C^1$-curve, \cref{def:length_c1_curve,def:length_c0_curve} agree. Notice that
\[\lvert\gamma'(t)\rvert = \sqrt{D\gamma(t)\transp \cdot D\gamma(t)}\]
is exactly the Gram determinant of $D\gamma(t)$ in the sense of
\Cref{def:volume_gram_determinant} --- the $m=1$ case of \cref{ex:gram_m1_length}, now
applied pointwise along the curve. This is why the length of a curve is the special case
$d=1$ of \cref{def:d_volume_parametrized_submanifold}: a curve is a $1$-dimensional
parametrized submanifold, and $L(\gamma) = \vol_1(\gamma([a,b]))$.
\end{remark}

% Supplement: Sascha Brack/Class Notes/Week_09_Notes_Friday.pdf, p. 4
\begin{example}[Line integral]
\label{ex:line_integral}
Let $\phi : [-\pi/2, \pi/2] \to \mathbb{R}^2$ be given by $\phi(\theta) = \begin{pmatrix} 3\sin\theta \\ 4\sin\theta \end{pmatrix}$.
Compute the length of the curve $M := \operatorname{Im}\phi$, i.e.\ $\int_M 1 \mathop{d\text{vol}}_1$.

First, we compute the Jacobian (which here is just the derivative vector):
\[\Jac\phi(\theta) = \begin{pmatrix} 3\cos\theta \\ 4\cos\theta \end{pmatrix}.\]
Then we compute the Gram determinant (which is the length of the tangent vector):
\[\sqrt{\det(\Jac\phi(\theta)\transp \Jac\phi(\theta))} = \sqrt{9\cos^2\theta + 16\cos^2\theta} = \sqrt{25\cos^2\theta} = 5|\cos\theta|.\]
Since $\cos\theta \geq 0$ for $\theta \in [-\pi/2, \pi/2]$, this simplifies to $5\cos\theta$.
The length is:
\[\int_{-\pi/2}^{\pi/2} 5\cos\theta \,d\theta = 5\big[\sin\theta\big]_{-\pi/2}^{\pi/2} = 5(1 - (-1)) = 10.\]
\end{example}

% Supplement: Sascha Brack/Class Notes/Week_09_Notes_Friday.pdf, p. 6
\begin{example}[Spiral length]
\label{ex:spiral_length}
Consider the spiral parametrized by $\phi : (0,2\pi) \to \mathbb{R}^3$,
\[\phi(t) = \begin{pmatrix} \cos(\omega t) \\ \sin(\omega t) \\ vt \end{pmatrix},\]
where $v \in \mathbb{R}\setminus\{0\}$ and $\omega \in (0,\infty)$. Compute the length of the spiral $\operatorname{Im}\phi$.

We compute the Jacobian:
\[\Jac\phi(t) = \begin{pmatrix} -\omega\sin(\omega t) \\ \omega\cos(\omega t) \\ v \end{pmatrix}.\]
The Gram determinant is:
\[\sqrt{\Jac\phi(t)\transp \Jac\phi(t)} = \sqrt{\omega^2\sin^2(\omega t) + \omega^2\cos^2(\omega t) + v^2} = \sqrt{\omega^2 + v^2}.\]
The length of the spiral is the integral over the parameter domain:
\[\int_0^{2\pi} \sqrt{\omega^2 + v^2} \,dt = 2\pi\sqrt{\omega^2 + v^2}.\]
\end{example}

% Supplement: Sascha Brack/Class Notes/Week_09_Notes_Monday.pdf, p. 8
\begin{example}[Line integral of a function on a circle]
\label{ex:line_integral_function_circle}
Given a function $f: \mathbb{R}^3 \to \mathbb{R}$, $f(x,y,z) = x^2+z^2$. We want to integrate this function on the circle $M := \{(x,y,z) \in \mathbb{R}^3 \mid x^2+y^2=1,\ z=2\}$.
We parametrize the circle by $\phi: (0,2\pi) \to \mathbb{R}^3$,
\[\phi(\varphi) = \begin{pmatrix} \cos\varphi \\ \sin\varphi \\ 2 \end{pmatrix}.\]
The Jacobian is $\Jac\phi(\varphi) = \phi'(\varphi) = \begin{pmatrix} -\sin\varphi \\ \cos\varphi \\ 0 \end{pmatrix}$.
The Gram determinant is $\sqrt{\Jac\phi\transp \Jac\phi} = |\phi'(\varphi)| = \sqrt{\sin^2\varphi + \cos^2\varphi} = 1$.
The function evaluated on the curve is $f(\phi(\varphi)) = \cos^2\varphi + 4$.
The line integral is thus:
\begin{align*}
\int_M f \mathop{d\text{vol}}_1 &= \int_0^{2\pi} f(\phi(\varphi)) |\phi'(\varphi)| \,d\varphi \\
&= \int_0^{2\pi} (\cos^2\varphi + 4) \cdot 1 \,d\varphi \\
&= \int_0^{2\pi} \cos^2\varphi \,d\varphi + \int_0^{2\pi} 4 \,d\varphi \\
&= \pi + 8\pi = 9\pi.
\end{align*}
\end{example}


% --- exercises relocated here from the dissolved sheet 9 ---

% Originally: content/week-09.tex, \exercisesheet{9}, problem 9.2
% Source: exercises/Ex9_Analysis2_eng.pdf

% Source: exercises/Ex9_Analysis2_eng.pdf, p. 1
\begin{exercise}[9.2 --- Multiple integrals]
\label{ex:9.2}
\begin{enumerate}[label=\textbf{\arabic*.}]
  \item Let $D := [0,2]\times[0,1]$. Calculate $\displaystyle\iint_D (x^3+3x^2y+y^3)\,dx\,dy$.
  \item Let $D \subset \mathbb{R}^2$ be the interior of the triangle with vertices $(0,0)$,
    $(0,\pi)$, and $(\pi,\pi)$. Calculate $\displaystyle\iint_D x\cos(x+y)\,dx\,dy$.
  \item Let $D := \{(x,y)\in\mathbb{R}^2 \mid x>1,\ y>1,\ x+y<3\}$. Calculate
    $\displaystyle\iint_D \frac{1}{(x+y)^3}\,dx\,dy$.
\end{enumerate}

% Source: Corsin Nick/Class Notes/Week 9.pdf, p. 1
\textit{Corsin's hint (part 2):} you can parametrize the triangle as
\[D = \{(x,y)\in\mathbb{R}^2 : 0 \leq x \leq y \leq \pi\}.\]
\end{exercise}
